\section{Data and experimental design}\label{sec:data}

\subsection{Sources, panels and snapshots}\label{sec:panels}

All series are monthly. Consumer price indices (all items) come from the IMF, through DBnomics for the first snapshot and the IMF SDMX API for the later ones. The exchange rate is the monthly average of domestic currency per US dollar (IMF IFS), and the Brent crude price is from the IMF primary commodity price series; both enter the LightGBM features. Dated snapshots with SHA-256 hashes were taken (two for the design countries and one file per hold-out panel; the hashes are in the repository).

\paragraph{Design countries (13).} Brazil, Mexico, Chile, Colombia, Peru, South Africa, Indonesia (core long-history group); Kazakhstan, Russia, Azerbaijan, Uzbekistan (regional group); Nigeria, Mongolia (control group). These countries shaped the design (the start date, the contrasts and the horizons). No export-share rule was applied to select them (Section~\ref{sec:scope}).

\paragraph{First hold-out panel (7).} Bolivia, Paraguay, Uruguay, Ghana, Georgia, Kyrgyzstan, Angola. The candidate list (14 names in a fixed order) and the eligibility rule (at least 220 months of IMF monthly CPI and exchange rate within 2000-01 to 2026-06, at most 3 missing CPI months, first 10 eligible candidates kept) were written before any forecast was computed for these countries. Only 7 candidates met the rule, and it was not relaxed. Target months run from 2009-02 to 2026-06.

\paragraph{Second hold-out panel (10, extension).} Zambia, Armenia, Tanzania, Uganda, Kenya, Iraq, Botswana, Namibia, Mozambique, C\^ote d'Ivoire. The candidate list and the rule (at least 170 months within 2010-01 to 2026-06, at most 3 missing CPI months, first 10 eligible candidates) were written before any forecast of any kind was computed for these countries. Each series has 198 months from 2010-01; six of the countries failed the 220-month requirement of the first rule, and no country is in two panels. After the 60-month training minimum and a warm-up of 48 there are 89 pooled calendar months at $h=3$ and 83 at $h=6$, and after a warm-up of 96 there are 41 and 35.

\paragraph{Data handling.} The two design snapshots overlap from 2024-01 to 2025-05 and agree except for constant rebasing of the index level for Kazakhstan and Nigeria, small revisions for Russia (up to 0.07 percent) and Mongolia (up to 0.05 percent), and a methodology break in the Nigerian series at 2025-01. The loader chains the snapshots on the CPI level at their last common month, so growth rates are unaffected (an early version without rescaling produced a false 74 percent fall in the Kazakhstani index, fixed before the run and now covered by a regression test). Nigerian months from 2025-01 are not appended, and Russia is cut into stretches of at least 60 consecutive months.

\subsection{Exploratory and confirmatory samples}\label{sec:samples}

Target months before 2025-07 are labelled exploratory. Target months from 2025-07 on the 13 design countries are labelled confirmatory in the code, but with at most about 16 months per country they cannot carry a claim and the run is degenerate (Section~\ref{sec:other}). The main confirmatory pool became the 7 first hold-out countries, and the extension adds the 10 countries of the second panel. Exploratory and hold-out results are kept separate throughout, and no threshold or window was tuned on hold-out data or to raise a result. The nominal level is 0.90, the primary warm-up is $n_{\mathrm{cal}}=48$ residuals and a sensitivity run uses $n_{\mathrm{cal}}=96$ for all methods; the latter is not promoted to the primary configuration.

\subsection{What ``pre-specified'' means and the sequence of events}\label{sec:prespec}

``Pre-specified'' means that a rule was committed to the repository before the run it governs. It does not mean registration in an external registry: the author is a single person, all entries are dated 2026-11-18 or 2026-11-19, and the order is attested only by a commit history kept by the same person. The decision rules R1--R5 were written after the exploratory results at both warm-ups were known. The sequence was as follows (commit hashes are in the repository change log).
\begin{enumerate}
\item The shock rule was frozen before any forecast. A committee review scored it 5.5 out of 10 (an exchange-rate flag measures currency turbulence, not an inflation shock), and its validity checks, run afterwards, gave a lift of 1.49 to 1.87 against the committee bar of 2 and a weak rank correlation with forecast error (0.075 and 0.050); stress is therefore secondary and exploratory, and the main contribution does not depend on shocks.
\item Protocol draft 1 was written before any forecast or interval was computed and revised after a committee review that scored it 5.5 out of 10; the third entry fixed the engine decisions (start 2000-01, warm-up 48, window 120, step 0.02) before results existed, and the first run was exploratory only (target months to 2025-06).
\item The fourth entry made amendments after the exploratory result and before any hold-out forecast: calibration error replaced pooled coverage as the primary outcome, the reference changed from split to rolling conformal, DRACP against ACI became an equivalence question with a 2~pp margin, warm-up 96 was added as a sensitivity run, the 7 hold-out countries became the main confirmatory pool and block length became at least 12. The entry attributes them to logical defects of the design and not to significance, but they were made after the exploratory results were known.
\item The fifth entry specified the hold-out runs (A: warm-up 48, primary; B: warm-up 96, sensitivity; C: an instant-feedback diagnostic at $h=12$) and the rules R1--R5 before any forecast for the 7 countries, and states that no rule is changed after the result is seen. The hold-out run was the first and only look at these countries.
\item The sixth entry (robustness analyses S1--S5) was written after the hold-out result and before they were computed; the seventh (post hoc S6, Kazakhstan alone, and S7, native intervals) was requested by a review of the manuscript and recorded before computing; both are secondary and do not replace R1--R5.
\item The eighth entry specified the extension (second panel, Chronos, rules E1--E5) after R1--R5 had been applied and after the R2 outcome on the first hold-out panel and the robustness results were known; it changes no earlier rule or result.
\end{enumerate}

\subsection{Procedural boundary}\label{sec:boundary}

The hold-out countries are untouched with respect to countries only. They are independent of the design sample by country but share calendar months with it (the global shocks of 2020 to 2022 fall in both samples). The methods, contrasts, horizons, the 2~pp margin and the 3~pp and 7~pp thresholds were chosen on the 13 design countries, and the author had seen their post-2025-06 months before the amendments. The hold-out evaluation is therefore a test on new countries and not a fully blind test of the design. The second panel is independent of both earlier samples by country, shares calendar months with them (2010 to 2026), and its rules E1--E5 were written by an author who knew every earlier result, including the R2 outcome on the first panel; it is a replication on new countries and not a blind test.

\subsection{Decision rules R1--R5 and extension rules E1--E5}\label{sec:rules}

The rules R1--R5 apply to LightGBM, the primary family with Holm correction over four tests, pooled over the 7 first hold-out countries with block length 12 (Table~\ref{tab:rules}). \textbf{R1:} ACI is better calibrated than split if ACI minus split is negative with Holm $p<0.05$ at both $h=3$ and $h=6$. \textbf{R2:} ``ACI improves on rolling'' only if ACI minus rolling is negative with Holm $p<0.05$ at the horizon, otherwise ``not shown''. \textbf{R3:} ``no advantage, equivalent within 2~pp'' only if equivalence holds at both warm-ups; ``DRACP better'' only if DRACP minus ACI is negative with Holm $p<0.05$ at both; otherwise ``inconclusive''. \textbf{R4:} delayed feedback explains the $h=12$ shortfall if the coverage of ACI with instant feedback exceeds the delayed one by at least 3~pp. \textbf{R5:} a contrast with an undefined MDE or an MDE above 7~pp is inconclusive regardless of its $p$-value.

\begin{table}[t]
\centering
\caption{Decision rules R1--R5 and extension rules E1--E5 as pre-specified, with their recorded outcomes.}
\label{tab:rules}
\small
\fittable{\begin{tabular}{p{0.05\linewidth}p{0.36\linewidth}p{0.11\linewidth}p{0.38\linewidth}}
\toprule
Rule & Statement & Status & Outcome as recorded \\
\midrule
\multicolumn{4}{l}{\emph{First hold-out panel (7 countries, LightGBM, Holm over four tests)}} \\
R1 & ACI better calibrated than split: ACI $-$ split negative with Holm $p<0.05$ at $h=3$ and $h=6$ & Not met & Not met in the primary run (warm-up 48); met at warm-up 96 (sensitivity run, not promoted) \\
R2 & ACI improves on rolling: ACI $-$ rolling negative with Holm $p<0.05$ & Not shown & Not shown (MDE undefined, inconclusive); Holm $p$ 0.35 and 1.0 at warm-up 48 \\
R3 & DRACP versus ACI: equivalence within 2 pp at both warm-ups, or DRACP better with Holm $p<0.05$ at both & Inconclusive & Equivalent at warm-up 48, not shown at 96; the verdict is not stable \\
R4 & Delayed feedback explains the $h=12$ shortfall: ACI coverage with instant feedback higher by at least 3 pp & Not met & Not met: +2.4 pp against the 3 pp threshold (coverage 0.884 delayed, 0.908 instant) \\
R5 & Power: MDE undefined or above 7 pp is inconclusive regardless of $p$ & Inconclusive & All four primary contrasts inconclusive (MDE undefined on the full hold-out sample) \\
\midrule
\multicolumn{4}{l}{\emph{Extension, second hold-out panel (10 countries, four base models, warm-up 48 primary)}} \\
E1 & Replication of R1, R2, R3 on new countries, same rule and direction & Partly & R1 met for all four base models at both horizons (not met at warm-up 96); R2 met at $h=3$ for LightGBM, ARIMA and naive B and at $h=6$ for ARIMA and naive B (Chronos at $h=3$ inconclusive under E5); R3 inconclusive \\
E2 & Dependence on the base model: contrasts side by side, no threshold & Descriptive & R1 and R3 do not depend on the base model; R2 at $h=6$ does (ARIMA, naive B met; LightGBM, Chronos not met) \\
E3 & Point accuracy of Chronos: relative MAE with bootstrap interval, descriptive & Descriptive & Relative MAE of Chronos on clean months 0.774, 0.699, 0.654, 0.895 at $h=1,3,6,12$; no interval favours LightGBM or ARIMA \\
E4 & Leakage: accuracy claims only from target months from 2025-01 & Applied & 22 months, 567 country-month cells per horizon, 30 countries; the 2025-01 cutoff is an assumption, not verified \\
E5 & Power: MDE undefined or above 7 pp is inconclusive & Applied & MDE of ACI $-$ rolling 1.2 to 6.6 pp at warm-up 48, undefined for Chronos at $h=3$; undefined for every DRACP $-$ ACI contrast and for all contrasts at warm-up 96 \\
\bottomrule
\end{tabular}
}
\end{table}

The extension (eighth entry, 2026-11-19) runs all four base models (LightGBM, ARIMA, naive B, Chronos) on the second panel at warm-up 48 (primary) and 96 (sensitivity, not promoted), starting at 2010-01, and re-runs the earlier 20 countries from 2010-01 as a secondary analysis, whose numbers are therefore not comparable cell by cell with the earlier sections. Its rules are applied separately for each base model. \textbf{E1:} a statement of R1, R2 or R3 is ``replicated'' only if the same rule is met in the same direction; warm-up 96 results are a sensitivity analysis and are not promoted. \textbf{E2:} contrasts for the four base models are reported side by side, without a threshold. \textbf{E3:} Chronos accuracy is descriptive. \textbf{E4:} accuracy claims for Chronos use only target months from 2025-01 pooled over all 30 countries, with the thinness of the sample stated, while calibration contrasts use all months with the flag. \textbf{E5:} a contrast with an undefined MDE or an MDE above 7~pp is inconclusive. Holm is applied within each base model and not across them; the four base models share countries and months, so they are not independent replications, and no statement rests on pooling them. Two deviations from the protocol are stated: calibration contrasts for the earlier 20 countries were analysed only for LightGBM and Chronos although four base models were run, and the all-months accuracy table has no intervals and no ETS, which was not part of the extension runs. A first analysis of the new panel failed because the country code NA (Namibia) was read as a missing value; the reader was fixed, the forecasts were not affected, and only the analysis was repeated.

\subsection{Commodity dependence of the panel}\label{sec:scope}

``Commodity-dependent'' is defined by a measured criterion: commodities at least 60\% of merchandise exports, averaged over 2021--2023 (UNCTAD, \emph{The State of Commodity Dependence 2025}). It was applied after the panels were chosen (post hoc), and part of the shares had been read before the criterion was written down. Table~\ref{tab:commodity} shows that 27 of the 30 economies qualify; Mexico, Indonesia and Georgia do not and are kept as comparison cases, and Uruguay (82.9\%) is strongly dependent. The hold-out sets followed IMF data availability under fixed rules. No comparison of dependent with other economies is possible, and the descriptive calibration error by class is inconclusive.

\begin{table}[t]
\centering
\caption{Commodity dependence of the three panels (UNCTAD, share of commodities in merchandise exports, average 2021--2023; classification applied post hoc).}
\label{tab:commodity}
\small
\fittable{\begin{tabular}{lrrrrl}
\toprule
Panel & Countries & $\geq 60\%$ & 40--60\% & $<40\%$ & Not qualifying at 60\% (share, \%) \\
\midrule
Design & 13 & 11 & 1 & 1 & Mexico (18.2), Indonesia (55.9) \\
Hold-out 1 & 7 & 6 & 1 & 0 & Georgia (46.4) \\
Hold-out 2 & 10 & 10 & 0 & 0 & none \\
All & 30 & 27 & 2 & 1 & \\
\bottomrule
\end{tabular}}
\end{table}

\subsection{Pipeline and platform (engineering)}\label{sec:architecture}

The figures and tables of the paper are generated from committed result files by a single script; numbers in the running text were written from the project manuscript and result documents. A platform module serves calibrated inflation intervals for the 13 design economies through a FastAPI endpoint and a React page that shows the measured backtest coverage next to the nominal level. It forms the point forecast as $h$ times an exponentially weighted moving average of past monthly log growth, issues intervals by ACI over a window of 120 months with step 0.02 and delayed feedback, and requires at least 30 revealed residuals (against 48 in the benchmark). The benchmark results bear on the module only through the interval method; ACI was chosen as an adaptive method that is better calibrated than one-shot calibration and native intervals and is simple to explain, not because it was shown to be better than a rolling window on all panels, and the module is an engineering application and not a scientific contribution. A service built on the endpoint answers questions about a country and horizon in English, Russian or Kazakh; its evaluation was specified in the eighth entry before any request was made. \ExplainerStatus{}
